\chapter*{General Introduction}
\addcontentsline{toc}{chapter}{General Introduction}
\markboth{General Introduction}{}
\label{ch:intro}
\lofgroup{General Introduction}

\section*{Context}

Diabetes is one of the fastest-growing chronic diseases in the world: it affects an
estimated 537 million adults, and this number is expected to reach 783 million by
2045~\cite{idf2021atlas}. Its most frequent complication in the eye, diabetic retinopathy
(DR), affects about 103 million people and is a leading cause of preventable blindness in
adults of working age~\cite{teo2021global,cheung2010dr}. The tragedy of DR is that it is
silent: vision is usually preserved until the disease is advanced, whereas treatment is
most effective in the early stages. Regular screening of the retina of every diabetic
patient is therefore the only way to prevent blindness at the scale of a population.

Screening relies on fundus photographs graded by trained experts on the five-level
International Clinical Diabetic Retinopathy (ICDR) scale~\cite{wilkinson2003icdr}. The number of images to grade now far
exceeds the number of available graders, especially in countries where ophthalmologists
are scarce. Automated grading by deep learning has become the natural answer, and in recent years
research has raised the agreement between algorithms and experts from a quadratic
weighted kappa (QWK) of $0.840$ (EfficientNet-B3, 2019~\cite{tan2019efficientnet}) to
$0.920$ (RETFound, 2023~\cite{zhou2023retfound}).

\section*{Problem Statement}

Despite this progress, no system reviewed in this thesis can be used, as it stands, in a
national screening programme. The reason is not a lack of data or of computing power; it
is architectural. Three limitations recur in every published system:
\begin{itemize}[leftmargin=1.4em,itemsep=3pt]
\item \textbf{L1 -- no relational context.} Each fundus image is graded in isolation, so
patients who look alike, and who often share a grade, never inform one another.
\item \textbf{L2 -- no vascular topology.} Convolutional and transformer networks see local
texture; they cannot measure the global structure of the vascular tree (its connected
components and loops), which changes as the disease progresses.
\item \textbf{L3 -- no right to abstain.} Every image receives a grade, even when it is
ambiguous. The most difficult cases, at the boundary between moderate and severe disease,
are forced through, and the rate of severe cases that are not referred (Sev-NR) stays
above the $2\%$ safety threshold used in this thesis.
\end{itemize}
These limitations are coupled: a graph needs topological information to be useful,
topology alone gives only a partial gain, and abstention only reaches the safety threshold
if the underlying predictions are accurate enough. They must therefore be addressed
together, in a single architecture. \Cref{fig:map} summarises how each limitation leads to
a research question, to a component of the proposed system and to the result obtained in
this thesis.

% Figure 2 (General Introduction) --- map of the argument
\begin{figure}[H]
\centering
\resizebox{0.9\linewidth}{!}{%
\begin{tikzpicture}[
  font=\footnotesize,
  box/.style={draw, rounded corners=2pt, align=center, inner sep=3pt,
              minimum height=1.75cm, text width=3.6cm},
  hd/.style={font=\small\bfseries, text=navy},
  arr/.style={-{Latex[length=1.8mm]}, thick, black!55}
]
\def\dx{4.05}\def\dy{2.3}
\node[hd, text=limred] at (0,0.3) {Structural limitation};
\node[hd] at (\dx,0.3) {Research question};
\node[hd] at (2*\dx,0.3) {What DOTS adds};
\node[hd, text=okgreen] at (3*\dx,0.3) {What was shown};

\node[box, fill=failf, draw=failred] (l1) at (0,-1.3) {\textbf{L1 -- no graph}\\Each image is scored alone;\\neighbours are ignored};
\node[box, fill=black!4, draw=black!50] (q1) at (\dx,-1.3) {Does a similarity graph add information that a CNN cannot recover alone?};
\node[box, fill=axisRf, draw=axisR] (s1) at (2*\dx,-1.3) {\textbf{Axis R}\\Dual-similarity graph\\GraphSAGE, $k{=}15$};
\node[box, fill=okf, draw=okgreen] (e1) at (3*\dx,-1.3) {\textbf{H1 confirmed}\\$t_4{=}3.82$, $p{=}0.009$\\$d_z{=}1.71$; ${+}4.0$ pp QWK};

\node[box, fill=axisSf, draw=axisS] (l2) at (0,-1.3-\dy) {\textbf{L2 -- no topology}\\CNNs see local texture,\\not vessel loops ($H_1$)};
\node[box, fill=black!4, draw=black!50] (q2) at (\dx,-1.3-\dy) {Do $H_0/H_1$ persistence features track ICDR severity?};
\node[box, fill=axisSf, draw=axisS] (s2) at (2*\dx,-1.3-\dy) {\textbf{Axis S}\\7-D Vietoris--Rips\\vector (GUDHI)};
\node[box, fill=okf, draw=okgreen] (e2) at (3*\dx,-1.3-\dy) {\textbf{H2 confirmed}\\$F(4,4995){=}1308.9$\\$\eta^2{=}0.51$ (large)};

\node[box, fill=axisOf, draw=axisO] (l3) at (0,-1.3-2*\dy) {\textbf{L3 -- no safety}\\Every image is graded;\\Sev-NR $>2\%$ before DOTS};
\node[box, fill=black!4, draw=black!50] (q3) at (\dx,-1.3-2*\dy) {Do the three axes meet the seven criteria, and does any subset fail?};
\node[box, fill=axisOf, draw=axisO] (s3) at (2*\dx,-1.3-2*\dy) {\textbf{Axis O}\\RETFound $+$ CORAL\\$+$ entropy abstention};
\node[box, fill=okf, draw=okgreen] (e3) at (3*\dx,-1.3-2*\dy) {\textbf{H3a / H3b}\\$7/7$ at point estimate\\$6/6$ subsets fail};

\foreach \a/\b in {l1/q1, q1/s1, s1/e1, l2/q2, q2/s2, s2/e2, l3/q3, q3/s3, s3/e3}
  \draw[arr] (\a) -- (\b);

\node[draw=okgreen, fill=okf, rounded corners=2pt, inner sep=5pt, font=\footnotesize,
      minimum width=15.75cm, align=center] at (1.5*\dx,-1.3-3*\dy+0.35)
  {\textbf{DOTS (all three axes):}\enspace QWK $=0.951$\enspace$\cdot$\enspace Sev-NR $=1.9\%$ $[1.4,\,2.4]$\enspace$\cdot$\enspace AUC $=0.988$\enspace$\cdot$\enspace $0\%$ CORAL violations\\
   H3b: every proper subset misses the $2\%$ safety rule; only the full system passes at the point estimate.};
\end{tikzpicture}}
\caption[Map of the thesis argument: limitations, components and results]{Map of
the argument. Each row links one structural limitation (L1--L3) to its research
question, the DOTS component that addresses it, and the result of the pre-specified
test. The bottom bar gives the full three-axis system.}
\label{fig:map}
\end{figure}


\section*{Objectives and Value of the Work}

This thesis aims to design, validate and engineer a DR grading system that resolves the
three limitations at once. More precisely, it aims to:
\begin{itemize}[leftmargin=1.4em,itemsep=3pt]
\item build an inductive population graph that lets similar patients share information,
while grading a new patient in constant time;
\item describe the retinal vascular network with persistent-homology features and use them
both as patient features and inside the similarity that builds the graph;
\item add ordinal prediction with guaranteed rank consistency and a calibrated abstention
mechanism, so that ambiguous cases are sent to a human grader;
\item validate each step with statistical tests fixed before the experiments, and engineer
the result as medical-device software compatible with current regulations.
\end{itemize}

The work thus pursues a double ambition. On the scientific side, it shows that the three
limitations can be resolved together and that, within the family of architectures studied,
none of the three components can be removed without losing the safety criterion. On the
practical side, it delivers a system that is accurate (QWK $=0.951$), safe at the point
estimate (Sev-NR $=1.9\%$), fast enough for national screening ($31.4$~ms per image), and
designed from the start to meet the requirements of the IEC~62304 and ISO~14971 standards
and of the EU AI Act. For a screening programme of $500{,}000$ patients a year, the
reduction of the non-referral rate corresponds to about $600$ fewer missed
sight-threatening cases under the assumptions of \cref{sec:impact}.

\section*{Research Methodology}

The work follows a hypothesis-driven methodology. Three research hypotheses structure the
thesis: \textbf{H1}, a population graph adds information that a per-image classifier cannot
recover; \textbf{H2}, persistent-homology features of the vessel network are associated with
disease severity; and \textbf{H3}, the combination of the three axes is sufficient (H3a) and,
within the architectures studied, necessary (H3b) to meet seven predefined clinical
criteria. For each hypothesis, the statistical test, the decision threshold and the
effect-size measure were fixed before the corresponding experiment was run, so that a
negative result could not be explained away afterwards; the falsification criteria are
listed in \cref{tab:falsif}. The components were introduced progressively, from a generic
graph template validated outside medicine, to TOPORET, and finally to DOTS, so that the
contribution of each component can be measured separately. All experiments use public
datasets, fixed random seeds and released code (\cref{app:repro}), so that they can be
reproduced independently.

\Cref{fig:pipe} gives an overview of the system obtained at the end of this process. Each
fundus image goes through six typed stages, and each stage belongs to one of the three
axes that address the limitations L1--L3.

% Figure (General Introduction) --- six-stage framework of the thesis
\begin{figure}[!htbp]
\centering
\resizebox{\linewidth}{!}{%
\begin{tikzpicture}[
  font=\footnotesize,
  box/.style={draw, rounded corners=2pt, align=center, inner sep=2.5pt,
              minimum height=1.9cm, text width=1.72cm},
  arr/.style={-{Latex[length=2.2mm]}, thick}
]
\node[box, fill=black!6, draw=black!55] (s1) {\textbf{S1}\\Pre-proc.\\CLAHE,\\skeleton};
\node[box, fill=black!6, draw=black!55, right=2.6mm of s1] (s2) {\textbf{S2}\\Backbone\\EfficientNet};
\node[box, fill=axisSf, draw=axisS, right=2.6mm of s2] (s3) {\textbf{S3}\\Persistent\\homology};
\node[box, fill=fusionf, draw=fusion, right=2.6mm of s3] (s4) {\textbf{S4}\\Fusion};
\node[box, fill=axisRf, draw=axisR, right=2.6mm of s4] (s5) {\textbf{S5}\\Population\\graph};
\node[box, fill=axisOf, draw=axisO, right=2.6mm of s5] (s6) {\textbf{S6}\\GNN and\\grading head};
\foreach \a/\b in {s1/s2,s2/s3,s3/s4,s4/s5,s5/s6} \draw[arr] (\a) -- (\b);
\node[left=2.5mm of s1, font=\footnotesize, align=center] (in) {Image\\$I_i$};
\node[right=2.5mm of s6, font=\footnotesize, align=center] (out) {Grade\\$\hat y_i$};
\draw[arr] (in) -- (s1); \draw[arr] (s6) -- (out);
\node[below=9mm of s3.south east, anchor=north, font=\footnotesize] (leg) {%
  \tikz\node[fill=axisRf,draw=axisR,minimum size=7pt,inner sep=0]{}; Axis R (L1)\quad
  \tikz\node[fill=axisSf,draw=axisS,minimum size=7pt,inner sep=0]{}; Axis S (L2)\quad
  \tikz\node[fill=axisOf,draw=axisO,minimum size=7pt,inner sep=0]{}; Axis O (L3)\quad
  \tikz\node[fill=fusionf,draw=fusion,minimum size=7pt,inner sep=0]{}; shared infrastructure};
\end{tikzpicture}}
\caption[Overview of the six-stage framework and the three axes]{Overview of the
six-stage framework with typed interfaces used in this thesis~\cite{belhadj2026enase}, and
the axis to which each stage belongs. TOPORET (\cref{ch:toporet}) instantiates all six
stages with vessel topology; in \cref{ch:dots}, DGTS replaces S3 by a topological
regulariser on the graph, and DOTS replaces S5--S6 by a feature refinement encoder and a
CORAL ordinal head (\cref{fig:dots}).}
\label{fig:pipe}
\end{figure}


The typed interfaces between stages make each stage testable on its own, a property that is exploited for software verification in \cref{ch:reg}.

\section*{Contributions}

The main contributions of this thesis, supported by eleven publications (see the List of
Publications), are the following.
\begin{description}[leftmargin=2.2em,style=sameline,itemsep=3pt]
\item[C1.] A GraphSAGE template with structural descriptors, validated on six
heterogeneous graph benchmarks outside medicine
\cite{belhadj2025icaart,belhadj2026invariants,belhadj2026lnaiinv}.
\item[C2.] Evidence, obtained with three independent methods, that topological and
structural features of the retina carry information that deep visual features miss,
including SAMF-GB, a CPU model more accurate than a GPU EfficientNet-B3 at about $1/13$ of
its cost \cite{belhadj2026hybrid,mezghich2026lnaitopo,belhadj2026icpram,belhadj2026samfgb}.
\item[C3.] TOPORET, the first inductive dual-similarity population graph for DR grading,
with statistical confirmation of the value of the graph (H1) and of the topological
features (H2) \cite{belhadj2025plos,belhadj2026enase,belhadj2026miua}.
\item[C4.] DOTS, to the best of our knowledge the first DR grading system that meets seven
predefined clinical-oriented criteria simultaneously on retrospective data (H3a), together
with an exhaustive proof that no reduced version of it does (H3b)
\cite{belhadj2026dots}.
\item[C5.] A software architecture and a regulatory mapping that make DOTS verifiable,
traceable and explainable by construction.
\end{description}

\Cref{tab:pubmap} shows which publications support each chapter of the thesis.

\begin{table}[H]
\centering
\caption{Publications supporting each chapter.}
\label{tab:pubmap}
\small
\begin{tabular}{@{}c P{6.4cm} P{4.2cm} l@{}}
\toprule
Ch. & Content & Contribution & Publications \\
\midrule
1 & Clinical problem, state of the art, limitations L1--L3 & Problem statement & -- \\
2 & Graphs, topology, ordinal prediction, metrics & Methodological background & -- \\
3 & Graph template outside medicine; structural evidence on DR & C1, C2 & P1--P4, B1, B2, J2 \\
4 & TOPORET population graph; tests of H1 and H2 & C3 & P5, P6, P7 \\
5 & DOTS system; sufficiency (H3a) and necessity (H3b) & C4 & J1 \\
6 & Software architecture and regulatory mapping & C5 & P6 \\
\bottomrule
\end{tabular}
\end{table}

All eleven publications are listed, with the contribution of each author, after the Bibliography.

\section*{Thesis Organisation}

The thesis is organised in six chapters, followed by a General
Conclusion that summarises the contributions, their limits and the perspectives they open:
\begin{enumerate}[leftmargin=1.6em,itemsep=3pt]
\item \textbf{Diabetic Retinopathy and Automated Grading} (\cref{ch:sota}): the disease,
its grading scale, screening programmes, two decades of methods for automated grading, and the limits
they share.
\item \textbf{Learning on Graphs, Topology and Ordinal Prediction}
(\cref{ch:background}): the methodological background -- graph neural networks,
persistent homology, ordinal regression, calibration and abstention -- and the evaluation
metrics.
\item \textbf{Structural and Topological Priors for Retinal Image Grading}
(\cref{ch:found}): validation of the graph template outside medicine, and first evidence
that structural features help DR grading.
\item \textbf{Topology-Aware Population Graph for Diabetic Retinopathy Grading}
(\cref{ch:toporet}): TOPORET, the dual-similarity population graph, and the tests of H1
and H2.
\item \textbf{Safe Ordinal Grading of Diabetic Retinopathy with Abstention} (\cref{ch:dots}): DOTS, the complete
system, its sufficiency (H3a) and necessity (H3b), its clinical impact and its comparison with the
literature.
\item \textbf{Software Architecture and Regulatory Compliance} (\cref{ch:reg}): how DOTS
is engineered as medical-device software.
\end{enumerate}

Readers with different backgrounds may follow different paths through the document.
Clinicians and regulatory readers will find a self-contained account in \cref{ch:sota},
\cref{ch:dots} and \cref{ch:reg}, together with the General Conclusion. Machine-learning
researchers will find the methodological background in \cref{ch:background} and the
technical contributions, with their algorithms and experiments, in
\cref{ch:found,ch:toporet,ch:dots}. Readers who wish to check the statistics or reproduce
the experiments will find the formal tests in \cref{app:stats} and the software
environment, seeds and computational cost in \cref{app:repro}.

